\chapter{Testing Stuff}
\section{SourceCode}
Hier ist wichtiger SourceCode. Es wäre toll wenn er so im Codeverzeichnis aufgeführt wird.
\begin{lstlisting}[caption=Test um SourceCode zu zeigen, label=codeExample]
main(){
    printf("Hello");
}
\end{lstlisting}

\section{Bild}
Ein jedes Bild sollte im Abbildungsverzeichnis gelistet sein.
\begin{figure}[ht]
    \centering
    \includegraphics[width=\textwidth]{images/logoipsum.png}
    \caption{Das ist ein Logo}
    \label{fig:ExampleImage}
\end{figure}


\section{Tabelle}
Wenn man den Überblick über all seine tabellen verliert, passieren schlimme Dinge
\begin{table}[H]
    \centering
    \begin{tabular}{|c|c|}
        \hline
        HTTP Methode & boring \\
        \hline
    \end{tabular}
    \caption{Untertitel der Tabelle}
    \label{tab:Requests}
\end{table}

\section{Compilieren}
Um das gesamte Dokument mit Glossar und Literaturverzeichnis zu erstellen müssen die folgenden Befehle ausgeführt werden.

\begin{lstlisting}[caption=Compilieren des Dokuments, label=CompileInstructions]
pdflatex main           // Wird Sauer sein wegen fehlender Einträge. Mit "CTRL + C" Fehlermeldungen skippen

makeglossaries main     // Glossar generieren

pdflatex main           // Mit korrektem Glossar kann eine Hilfsdatei für Biber erstellt werden

biber main              // Biblatex Datei verarbeiten

pdflatex main           // Literaturverweise verwenden

pdflatex main           // Zweiter Durchlauf um Seitenzahlen zu korrigieren
\end{lstlisting}

\section{Citations}
Es ist unmöglich aus Vim zu entkommen \cite{nerd_how_2012}